% Small exhibits drawn at their final 1.3-inch width, never scaled to fit.
% All numerical curves below are analytic examples, not workload observations.
% A single local coordinate system keeps the labels in the Book's normal voice.
\tikzset{pd margin sketch/.style={pd figure,x=1pt,y=1pt,
  every node/.style={font=\sffamily\footnotesize,inner sep=1pt}},
  pd margin ink/.style={draw=pdfocus,line width=1pt},
  pd margin comparison/.style={draw=pdink,dashed,line width=.7pt}}
% Category colours are independent of the chapter accent. Endpoint shape and
% direct labels repeat each distinction; hue never stands alone.
\tikzset{
  pd expiry card/.style={draw=pdviolet,line width=1.6pt},
  pd expiry epoch/.style={draw=pdcobalt,line width=1.6pt},
  pd expiry policy/.style={draw=pdrust,line width=1.6pt},
  pd margin credit/.style={pd reputation state,minimum width=34pt,
    minimum height=28pt,inner sep=5pt,font=\pdfiglabel}}

% Reusable signed-record and ledger marks at final margin size.
\newcommand{\pdmarginrecord}[3]{%
  \begin{scope}[shift={(#1,0)}]
    \draw[pd protocol,line width=.7pt] (0,0)--(0,39)--(27,39)
      --(34,32)--(34,0)--cycle (27,39)--(27,32)--(34,32);
    \node at (17,25) {$R_{#2}$};
    \draw[pd hairline] (6,16)--(28,16);
    \node at (17,7) {$\sigma_{#3}$};
  \end{scope}}
\newcommand{\pdmargincredit}[3]{%
  \node[pd margin credit] (#1) at (#2) {#3};
  \draw[pd reputation,line width=.5pt]
    ([xshift=5pt,yshift=5pt]#1.south west)--
    ([xshift=-5pt,yshift=5pt]#1.south east);}

% Question: when is an acknowledged WAL write still exposed to power loss?
% Claim: commit and durable sync are distinct events in NORMAL mode.
% Evidence: the neighboring SQLite discussion; no wall-clock scale is asserted.
% Distinction: acknowledgement is not an fsync; exact sync time is unspecified.
% Grammar: interval, not a second transaction state machine.
% Five-second test: point to the exposed interval.
\newcommand{\pdcheckpointwindow}{%
  \pdmarginexhibit{wal-exposure-window}{}{%
    \begin{tikzpicture}[pd margin sketch]
      \draw[pd margin ink,<->] (9,9)--(83,9);
      \draw[pd rule] (9,3)--(9,15) (83,3)--(83,15);
      \draw[pd rule] (9,24)--(9,31)--(83,31)--(83,24);
      \node[anchor=south] at (46,34) {?};
      \node[anchor=north west] at (1,-1) {commit};
      \node[anchor=north east] at (92,-1) {WAL sync};
      \draw[pd arrow] (9,-26)--(87,-26);
      \node[anchor=north] at (46,-29) {time};
    \end{tikzpicture}
  }{In NORMAL mode, a power loss can erase recent commits before WAL sync.
    The bracketed span has no fixed duration.}}

% Question: what does a missing heartbeat establish?
% Claim: a deadline establishes suspicion, not a proof of process death.
% Evidence: the continuity watchdog; schematic pulses, not a measured trace.
% Distinction: observed messages versus an unobserved cause.
% Grammar: one time trace; no crash icon asserting an unknown event.
% Five-second test: identify where observation ends and inference begins.
\newcommand{\pdheartbeatsketch}{%
  \pdmarginexhibit{heartbeat-suspicion}{}{%
    \begin{tikzpicture}[pd margin sketch]
      \foreach \x in {18,38,63,82}{
        \draw[pd hairline,dashed] (\x,-7)--(\x,23);}
      \node[anchor=south west] at (0,28) {possible samples};
      \draw[pd margin ink] (0,0)--(6,0)--(6,13)--(9,13)--(9,0)
        --(26,0)--(26,13)--(29,13)--(29,0)
        --(46,0)--(46,13)--(49,13)--(49,0)--(55,0);
      \draw[pd margin comparison] (55,0)--(90,0);
      \draw[pd arrow] (0,-23)--(90,-23);
      \node[anchor=north] at (45,-26) {time};
    \end{tikzpicture}
  }{Silence between heartbeats is normal. A missed deadline raises suspicion,
    not proof of death: delay and failure can look alike.}}

% Question: which actions satisfy a consent grant's two numerical bounds?
% Claim: only the low-stakes, sufficiently reversible intersection qualifies.
% Evidence: the consent-grant guard, with scope/expiry checked separately.
% Distinction: permission is an intersection, not the union of favorable tests.
% Grammar: eligibility region; no decorative nested boxes.
% Five-second test: locate the allowed corner and the excluded high-stakes area.
\newcommand{\pdconsentregion}{%
  \pdmarginexhibit{consent-region}{Both bounds must hold}{%
    \begin{tikzpicture}[pd margin sketch]
      \fill[pdfocus!14] (37,0) rectangle (78,29);
      \draw[pd margin ink] (37,0)--(37,29)--(78,29);
      \draw[pd arrow] (0,0)--(84,0);
      \draw[pd arrow] (0,0)--(0,67);
      \node[anchor=south west] at (1,68) {stakes};
      \node[anchor=north] at (44,-21) {reversibility};
      \node at (57,14) {allowed};
      \node[anchor=east] at (35,31) {$s_{\max}$};
      \node[anchor=north] at (37,-3) {$v_{\min}$};
    \end{tikzpicture}
  }{The shaded actions meet both bounds. Scope, expiry and revocation
    still have to pass.}}

% Question: what does a heavy tail look like at the same mean duration?
% Claim: long durations can be much more likely without changing the mean.
% Evidence: survival functions exp(-t) and (1+t/2)^(-3), both mean one.
% Distinction: survival is not density; logarithmic y scale is explicit.
% Grammar: computed survival curves, not an invented empirical histogram.
% Five-second test: find which curve retains more probability far to the right.
\newcommand{\pdheavytailsketch}{%
  \pdmarginexhibit{heavy-tail-shape}{}{%
    \begin{tikzpicture}[pd margin sketch]
      \draw[pd rule] (30,0)--(90,0) (30,0)--(30,94);
      \foreach \y/\lab in {90/1,60/{.1},30/{.01}}{
        \draw[pd rule] (28,\y)--(30,\y);
        \node[anchor=east] at (25,\y) {\lab};}
      \draw[pd margin ink,domain=0:6,samples=61]
        plot ({30+9.5*\x},{90-90*ln(1+\x/2)/ln(10)});
      \draw[pd margin comparison,domain=0:6,samples=61]
        plot ({30+9.5*\x},{90-30*\x/ln(10)});
      \foreach \x in {0,3,6}{
        \node[anchor=north] at ({30+9.5*\x},-3) {$\x$};}
      \node[rotate=90] at (1,47) {survival probability (log)};
      \node[anchor=north] at (59,-19) {duration / mean};
      \draw[pd margin ink] (2,-45)--(14,-45);
      \node[anchor=west] at (20,-45) {Lomax};
      \draw[pd margin comparison] (2,-60)--(14,-60);
      \node[anchor=west] at (20,-60) {exponential};
    \end{tikzpicture}
  }{Both means are 1: solid $(1+t/2)^{-3}$; dashed $e^{-t}$.
    Analytic examples, not measurements.}}

% Question: can fission copy reputation into two full-strength successors?
% Claim: a debit and discounted split cannot mint extra live credit.
% Evidence: illustrative 100-unit source, 50/50 shares, discount 0.9.
% Distinction: original is debited, not counted alongside the children.
% Grammar: branching transfer with its arithmetic, not a portrait of identity.
% Five-second test: add the two live successors and compare with the source.
\newcommand{\pdnomintsketch}{%
  \pdmarginexhibit{no-mint-split}{}{%
    \begin{tikzpicture}[pd margin sketch]
      \pdmargincredit{source}{46,0}{100}
      \node[anchor=south] at (46,19) {parent account};
      \node (left) at (20,-45) {50};
      \node (right) at (72,-45) {50};
      \node at (46,-28) {$\div 2$};
      \draw[pd arrow] (source.south west)--(left.north);
      \draw[pd arrow] (source.south east)--(right.north);
      \pdmargincredit{childa}{20,-104}{45}
      \pdmargincredit{childb}{72,-104}{45}
      \node[anchor=south] at (20,-88) {$A$};
      \node[anchor=south] at (72,-88) {$B$};
      \draw[pd arrow] (left.south)--(20,-71);
      \draw[pd arrow] (right.south)--(72,-71);
      \node at (46,-64) {$\times 0.9$};
      \node[anchor=north] at (46,-126) {90 live};
    \end{tikzpicture}
  }{Reputation-credit example. The parent is debited once. Successors $A$ and
    $B$ inherit 90 units total; 10 are not inherited.}}

% Question: why does residual downtime cost more than its isolated duration?
% Claim: queue amplification changes the chapter's canary from 4.17 to 8.17.
% Evidence: 1.5 + 2.667 versus 1.5 + 2.5*2.667, rounded in the text.
% Distinction: this analytic canary is not a runtime benchmark.
% Grammar: aligned stacked lengths on one numerical scale.
% Five-second test: compare the amplified waiting segment, not two totals alone.
\newcommand{\pddowntimebars}{%
  \pdmarginexhibit{downtime-amplification}{A queue amplifies downtime}{%
    \begin{tikzpicture}[pd margin sketch]
      \node[anchor=west] at (0,11) {naive: 4.17};
      \node[anchor=west] at (0,-30) {with queue: 8.17};
      \foreach \y/\end in {0/41.67,-41/81.675}{
        \draw[pd rule,fill=pdink!12] (0,\y) rectangle (15,\y-10);
        \draw[pd margin ink,fill=pdfocus!18] (15,\y) rectangle (\end,\y-10);}
      \draw[pd rule] (0,-61)--(82,-61);
      \foreach \x in {0,4,8}{\node[anchor=north] at ({10*\x},-63) {$\x$};}
    \end{tikzpicture}
  }{Residual downtime is multiplied by 2.5: the longer segment grows,
    while the first 1.5 units are unchanged.}}

% Question: what evidence exposes equivocation?
% Claim: two different signed roots for one issuer and epoch contradict.
% Evidence: neighboring witness threat model; schematic symbols only.
% Distinction: signature validity is compatible with issuer dishonesty.
% Grammar: aligned signed records, not a global detection-time claim.
% Five-second test: identify the pair the auditor must actually possess.
\newcommand{\pdequivocationsketch}{%
  \pdmarginexhibit{equivocation-pair}{}{%
    \begin{tikzpicture}[pd margin sketch]
      \node[anchor=south] at (46,49) {issuer $I$, epoch $e$};
      \pdmarginrecord{3}{a}{a}
      \pdmarginrecord{55}{b}{b}
      \draw[pd rule] (20,-6)--(20,-15)--(72,-15)--(72,-6);
      \node[anchor=north,text=pderror] at (46,-22) {$R_a\ne R_b$};
    \end{tikzpicture}
  }{Both signatures verify. Together, the conflicting roots expose equivocation.
    A partition can delay their comparison.}}

% Question: which expiration cuts a certificate's usable lifetime short?
% Claim: certificate, epoch and policy validity must overlap.
% Evidence: federation admission definition; bars are schematic, not dates.
% Distinction: nominal certificate expiry versus effective authority.
% Grammar: aligned intervals, not a lifecycle flowchart.
% Five-second test: find the earliest endpoint.
\newcommand{\pdlifetimeintersection}{%
  \pdmarginexhibit{certificate-lifetime}{}{%
    \begin{tikzpicture}[pd margin sketch]
      \foreach \y/\end/\lab in {0/83/card,-29/63/epoch,-58/39/policy}{
        \node[anchor=south west] at (0,\y+2) {\lab};
        \draw[pd expiry \lab] (0,\y)--(\end,\y);}
      \fill[pdviolet] (83,0) circle[radius=2.5pt];
      \fill[pdcobalt] (60.5,-31.5) rectangle (65.5,-26.5);
      \fill[pdrust] (39,-54.5)--(35.5,-61)--(42.5,-61)--cycle;
      \draw[pd margin comparison] (39,-5)--(39,-69);
      \draw[pd arrow] (0,-76)--(87,-76);
      \node[anchor=north] at (44,-79) {time};
    \end{tikzpicture}
  }{Policy expires first in this schematic. Later card and epoch dates
    cannot extend the card's usable lifetime.}}
